% Generated draft: requires mathematical and source review.
\documentclass[11pt]{article}
\usepackage[margin=25mm]{geometry}
\usepackage{fontspec}
\setmainfont{Latin Modern Roman}
\usepackage{amsmath,amssymb,mathtools,unicode-math}
\setmathfont{Latin Modern Math}
\usepackage{enumitem,needspace,xurl,hyperref}
\setlength{\parindent}{0pt}
\setlength{\parskip}{5pt}
\setlength{\emergencystretch}{3em}
\newcommand{\N}{\mathbb{N}}
\newcommand{\Z}{\mathbb{Z}}
\newcommand{\Q}{\mathbb{Q}}
\newcommand{\R}{\mathbb{R}}
\newcommand{\C}{\mathbb{C}}
\newcommand{\F}{\mathbb{F}}
\newcommand{\A}{\mathbb{A}}
\newcommand{\PP}{\mathbb{P}}
\newcommand{\E}{\mathbb{E}}
\newcommand{\eps}{\varepsilon}
\newcommand{\dd}{\,\mathrm{d}}
\begin{document}
\section*{UnsolvedMath ID: 6900006}
\textbf{Source problem code:} catalogue entry
\subsection{Definitions and mathematical statement}
For a finite graph \(G=(V,E)\), a placement \(p:V\to\R^d\) has a self-stress \(\omega:E\to\R\) when
\[
 \sum_{j:\{i,j\}\in E}\omega_{ij}(p_j-p_i)=0\qquad(i\in V).
\]
The source concerns strata in configuration spaces of tensegrities, but does not supply here its precise stratum equivalence or admissible placement domain.

In this projective geometry setting, ``Cayley algebra'' is understood as the Grassmann--Cayley calculus of joins and meets, not the octonion algebra. The question is to characterize when two systems of such geometric conditions define the same stratum. If \(D\) is the specified common admissible configuration domain and \(\mathcal C\) is a system of conditions, write
\[
 Z_D(\mathcal C)=\{p\in D:p\text{ satisfies every condition in }\mathcal C\}.
\]
Then the equivalence to be recognized is \(Z_D(\mathcal C_1)=Z_D(\mathcal C_2)\), with the source's additional stratum conventions retained. A generators-and-relations description should explain the transformations between descriptions representing the same stratum.

The supplied summary illustrates the problem using \(K_{3,3}\): the condition that six points lie on a conic is independent of their ordering and reportedly yields 60 Cayley algebra systems describing the same stratum. This is an illustration from the supplied summary, not a verified enumeration. The excerpt does not specify the coefficient field, degenerate configurations, inequations, or the precise allowable algebraic systems. Equality of unrestricted polynomial zero sets, equality of ideals, and equality of real configuration strata must not be assumed interchangeable.
\subsection{Short English statement}
Characterize when different Grassmann--Cayley systems of geometric conditions describe the same tensegrity configuration stratum, ideally through generators and relations among those descriptions. The admissible configurations and the exact stratum convention are missing from the short question.
\subsection{Sources}
\begin{itemize}
\item Local UnsolvedMath registry, permanent ID 6900006, source code catalogue entry; Oleg Karpenkov, 2020; update 2026-08-06.
\item Karpenkov, \emph{Open Problems on Configuration Spaces of Tensegrities} (2020), Problem 6, PDF page 5: \url{https://livrepository.liverpool.ac.uk/3022838/1/stratum.pdf}. Registry access date 2026-07-29; normalized PDF extraction; public institutional-repository author manuscript; current status and reuse terms need review.
\item The supplied summary gives the \(K_{3,3}\) conic example, the reported number 60, and the generators-and-relations motivation. It gives no additional reference establishing a complete characterization.
\end{itemize}
\end{document}
